	\begin{UseCase}{CU12.6}{Terminar Caso de uso}{
		Este caso de uso permite al actor terminar la edición de un caso de uso en el sistema, para que pueda ser revisado y en su caso liberado.
	}
		\UCitem{Actor}{\hyperlink{jefe}{Líder de Análisis}, \hyperlink{analista}{Analista}}
		\UCitem{Propósito}{Terminar la edición de un caso de uso para solicitar su revisión.}
		\UCitem{Entradas}{Ninguna}
		\UCitem{Salidas}{
			\begin{itemize}
				\item \cdtIdRef{MSG1}{Operación exitosa}: Se muestra en la pantalla \IUref{IU6}{Gestionar Casos de uso} para indicar que la operación fue exitosa.
				\item \cdtIdRef{MSG27}{Confirmación de termino}: Se muestra en la pantalla \IUref{IU6}{Gestionar Casos de uso} para que el actor confirme el término del caso de uso.
		\end{itemize}}
		\UCitem{Precondiciones}{
				Que el caso de uso se encuentre en estado ''Edición'' o ''Pendiente de corrección''.
		}
		\UCitem{Postcondiciones}{
				El caso de uso pasará a estado ''Revisión''.
		}
		\UCitem{Errores}{Ninguno}
		\UCitem{Tipo}{Secundario, extiende del caso de uso \UCref{CU12}{Gestionar Casos de uso}.}
	\end{UseCase}
%--------------------------------------
	\begin{UCtrayectoria}
		\UCpaso[\UCactor] Da clic en el icono \raisebox{-1mm}{\includegraphics[height=11pt]{images/Iconos/terminar}} del registro que desea realizar la operación de la pantalla \IUref{IU6}{Gestionar Casos de uso}.
		\UCpaso[\UCsist] Verifica que el estado del caso de uso sea ''Edición'' o ''Pendiente de corrección''. \hyperlink{CU12-5:TAA}{[Trayectoria A]}
		\UCpaso[\UCsist] Muestra el mensaje emergente \cdtIdRef{MSG27}{Confirmación de termino} con los botones \IUbutton{Aceptar} y \IUbutton{Cancelar} en la pantalla \IUref{IU6}{Gestionar Casos de uso}.
		\UCpaso[\UCactor] Oprime el botón \IUbutton{Aceptar}. \hyperlink{CU12-5:TAB}{[Trayectoria B]}
		\UCpaso[\UCsist] Cambia el estado del caso de uso a ''Revisión''.
		\UCpaso[\UCactor] Se muestra el mensaje \cdtIdRef{MSG1}{Operación exitosa} en la pantalla \IUref{IU6}{Gestionar Casos de uso}.
	\end{UCtrayectoria}		
%--------------------------------------
\hypertarget{CU12-5:TAA}{\textbf{Trayectoria alternativa A}}\\
\noindent \textbf{Condición:} El caso de uso que se desea terminar se encuentra en un estado diferente a ''Edición'' o ''Pendiente de corrección''.
\begin{enumerate}
	\UCpaso[\UCsist] Oculta el botón \raisebox{-1mm}{\includegraphics[height=11pt]{images/Iconos/terminar}} del caso que no se encuentra en estado de ''Edición'' o ''Pendiente de corrección''.
	\item[- -] - - {\em {Fin del caso de uso}}.
\end{enumerate}
%--------------------------------------
\hypertarget{CU12-5:TAB}{\textbf{Trayectoria alternativa B}}\\
\noindent \textbf{Condición:} El actor desea cancelar la operación.
\begin{enumerate}
	\UCpaso[\UCactor] Oprime el botón \IUbutton{Cancelar} de la pantalla emergente.
	\UCpaso[\UCsist] Muestra la pantalla \IUref{IU6}{Gestionar Casos de Uso}.
	\item[- -] - - {\em {Fin del caso de uso}}.%
\end{enumerate}